\SourceLeaf{219}%
\noindent holds by assumption, but that the proposition

\[ c\;[\text{1st figure}] = c\;[\text{2nd figure}] \]

\noindent would then have to be proved?

Given that the angles are arbitrary, this proof seems to me unnecessary.

These are the main outlines of the proof, and I await your verdict. In defense of my proof I add only that the second operation does indeed make the triangle $ABC$ vanish, but \textls{not} the angles of the triangle. However the lines lie, we always have

\[ J\hat{B}H = \beta, \qquad G\hat{C}F = \gamma, \qquad D\hat{A}E = \alpha \]

\noindent in the finite triangle as in the vanishing triangle, so that the sum

\[ J\hat{A}H + G\hat{A}F + D\hat{A}E \]

\noindent is always equal to the sum of the angles of a rectilinear triangle.

To prove the proposition for any triangle (whose angles are $A$, $B$, $C$), draw the lines $DG$, $EH$ so that

\[ \alpha = A, \]

\noindent and then take the angle $J\hat{A}H = B$, $G\hat{A}F = C$.

If $JAF$ is then a broken line $JAF'$ instead of a straight line, the angle $c$ will indeed have decreased by $dc$, and the angle $b$ will have increased by just as much. Their sum therefore remains unchanged, or we have what the rigor of the proof requires:

\[ b + c\;[\text{Fig. 1}] = b + c\;[\text{Fig. 2}].\} \]

\bigskip
\hrule
\bigskip

\Needspace{7\baselineskip}
\begin{center}
\textsc{Gauss} to \textsc{Schumacher}. Göttingen, 17 May 1831.
\end{center}

\mbox{.\ .\ .\ .\ .\ .\ .} On closer inspection, in what you write about parallel lines you have used an intermediate proposition in your syllogisms without stating it explicitly. It would have to read as follows:
